\section{Considerações Finais}

Este trabalho partiu da questão de se é possível combinar dados abertos para criar um sistema interativo que facilite a visualização das empresas de TI de Contagem.
A resposta é positiva, com ressalvas: a aplicação integra as bases da Receita Federal e do IBGE e permite consultar, filtrar e localizar 781 estabelecimentos, mas a coordenada é aproximada.

Os objetivos específicos foram atingidos, com ressalvas.
Na base única, a coordenada vem apenas do CEP, e a proporção de MEI não pôde ser afirmada com segurança; na aplicação, a agregação por área do município, prevista no início, não foi implementada.
Nos testes, a inspeção heurística apontou problemas de usabilidade a corrigir.

Como contribuição prática, o trabalho oferece uma consulta única e sem cadastro às empresas de TI do município, com aviso da precisão da coordenada.
Como contribuição acadêmica, documenta um procedimento reprodutível de integração de CNPJ, CNAE e CNEFE, com níveis de junção explícitos.
A principal lição foi que a ligação entre bases por endereço é a etapa mais frágil e que a taxa de acerto deve ser informada junto com o mapa.

\subsection{Limitações}

A CNAE é a atividade declarada no cadastro, e usar só a atividade principal deixa de fora estabelecimentos de outro ramo que também prestam serviço de TI.
A coordenada representa um endereço do CNEFE no mesmo CEP, e não a localização auditada do estabelecimento, e o campo de MEI não foi validado em outra fonte.
O banco guarda a fotografia de um mês, e a inspeção heurística, feita por três especialistas, não mede como usuários reais utilizam a aplicação.

\subsection{Trabalhos futuros}

Como continuidade, propõe-se incluir o tipo de logradouro na chave de junção, agrupar ou deslocar os marcadores sobrepostos no mapa e implementar a agregação por bairro, com o total e a densidade por quilômetro quadrado, a partir das malhas do IBGE.
Outras possibilidades são validar a condição de MEI em outra fonte, aplicar testes de usabilidade com usuários reais, automatizar a atualização mensal da base e estender a aplicação a outros municípios e setores.
